\documentclass[11pt]{article}
\usepackage[letterpaper,margin=1in]{geometry}
\usepackage[T1]{fontenc}
\usepackage[utf8]{inputenc}
\usepackage{lmodern,microtype}
\usepackage{amsmath,amssymb,mathtools,bm,slashed}
\usepackage{xcolor,enumitem,booktabs,array,needspace}
\usepackage[most]{tcolorbox}
\usepackage{fancyhdr,titlesec,parskip,hyperref}
\definecolor{qftblue}{HTML}{245A8D}
\definecolor{qftgreen}{HTML}{2F6B4F}
\definecolor{qftgray}{HTML}{F4F6F8}
\hypersetup{colorlinks=true,linkcolor=qftblue,urlcolor=qftblue,pdfauthor={Nima Lashkari}}
\pagestyle{fancy}\fancyhf{}
\lhead{PHYS 662: Quantum Field Theory I}\cfoot{\thepage}
\setlength{\headheight}{14pt}
\titleformat{\section}{\Large\bfseries\color{qftblue}}{\thesection.}{0.6em}{}
\titleformat{\subsection}{\large\bfseries\color{qftgreen}}{\thesubsection}{0.6em}{}
\setlist[enumerate]{leftmargin=2em,itemsep=0.25em,topsep=0.35em}
\setlist[itemize]{leftmargin=2em,itemsep=0.2em,topsep=0.35em}
\newcommand{\dd}{\mathrm d}
\newcommand{\ii}{\mathrm i}
\newcommand{\eps}{\varepsilon}
\newcommand{\tr}{\operatorname{tr}}
\newcommand{\Tr}{\operatorname{Tr}}
\newcommand{\diag}{\operatorname{diag}}
\newcommand{\Lag}{\mathcal L}
\newcommand{\D}{\mathcal D}
\newcommand{\vac}{\lvert0\rangle}
\newcommand{\vbra}{\langle0\rvert}
\newcommand{\Tord}{\mathrm T}
\newcommand{\ket}[1]{\lvert#1\rangle}
\newcommand{\bra}[1]{\langle#1\rvert}
\newtcolorbox{keyidea}[1][]{enhanced,breakable,colback=qftgray,colframe=qftblue,
 boxrule=0.7pt,arc=1.5mm,title={#1},fonttitle=\bfseries,coltitle=white}
\newtcolorbox{remarkbox}[1][]{enhanced,breakable,colback=white,colframe=qftgreen,
 boxrule=0.6pt,arc=1.5mm,title={#1},fonttitle=\bfseries,coltitle=white}
\newcommand{\intp}{\int\!\frac{\dd^4p}{(2\pi)^4}}
\newcommand{\MSbar}{\overline{\mathrm{MS}}}
\newcommand{\module}[1]{\clearpage\section{#1}}
\newcommand{\braket}[2]{\langle#1\mid#2\rangle}
\newcommand{\dPi}{\dd\Pi}
\newcommand{\cM}{\mathcal M}
\newcommand{\cH}{\mathcal H}
\newcommand{\cS}{\mathcal S}
\newcommand{\Zp}{Z_{\rm pole}}
\newcommand{\Kallen}{\mathcal K}
\newcommand{\lecturefront}[4]{%
\rhead{Lecture #1}\hypersetup{pdftitle={PHYS 662 Lecture #1: #2}}
\begin{document}\pagenumbering{gobble}
\begin{titlepage}\centering\vspace*{1.05in}
{\scshape\Large PHYS 662 --- Quantum Field Theory I\par}
\vspace{0.35in}{\Huge\bfseries\color{qftblue}Lecture #1\par}
\vspace{0.2in}{\LARGE\bfseries #2\par}
\vspace{0.2in}{\large\color{qftgreen}#3\par}
\vspace{0.6in}{\large Nima Lashkari\par}{\large Fall 2026\par}
\vfill\begin{minipage}{0.86\textwidth}\small #4\par\medskip
Course page: \url{https://nimalashkari.com/quantum-field-theory-2026-part-i/}
\end{minipage}\vspace{0.5in}\end{titlepage}
\pagenumbering{roman}\setcounter{tocdepth}{1}\tableofcontents
\vfill\begin{remarkbox}[Place in the course]
Lecture 26 constructs the physical photon Hilbert space and its propagators.
Lecture 27 develops QED amplitudes. Lecture 28 computes the photon and
fermion self-energies. Lecture 29 relates the vertex to charge
renormalization and the magnetic moment. Homework 9 closes this sequence.
\end{remarkbox}
\small\textbf{Reading and conventions.} The uploaded course slides are the
primary starting point. D.~Tong,
\href{https://www.damtp.cam.ac.uk/user/tong/qft/qfthtml/S6.html}{\textit{Quantum Field Theory}, Chapter 6},
is a companion for photon quantization and QED. For the vertex calculation,
see V.~Kaplunovsky,
\href{https://web2.ph.utexas.edu/~vadim/Classes/2020f/vertex.pdf}{\textit{QED Vertex Correction}},
especially the parameter integral for the Pauli form factor.
For further textbook study, see Peskin--Schroeder, \textit{An Introduction
to QFT}, Chapters 5--7 and 10.
All formulas here use the mostly-minus metric and propagators that include
\(\ii\). Dimensional regularization uses \(d=4-2\eps\), which differs from
some of the slides. The accompanying editorial memo records corrections
and the scope of the supplementary reading.
\normalsize\clearpage\pagenumbering{arabic}}
\lecturefront{29}{Loop Corrections to a Vertex}
{The Ward--Takahashi identity and the anomalous magnetic moment}
{The photon--electron vertex contains both a divergent correction related
to field normalization and a finite spin-dependent response. We separate
these pieces, derive the relation between charge and field
renormalization, and calculate the leading anomalous magnetic moment.
The final step explains why finite measured probabilities require
inclusive treatment of soft photons.}

\section{The vertex and the inverse propagator}
Let an electron of momentum \(p\) absorb a photon of momentum
\(q=p'-p\) and emerge with momentum \(p'\). Define the renormalized
amputated 1PI vertex by
\begin{equation}
 -\ii e\Gamma_R^\mu(p',p),\qquad
 \Gamma_R^\mu=\gamma^\mu+\Lambda^\mu_{\rm loop}+\delta_1\gamma^\mu.
\end{equation}
The factor \(-\ii e\) is removed in the definition of \(\Gamma_R\).
The loop correction \(\Lambda\) starts at order \(e^2\).
An amputated 1PI vertex is not by itself a scattering amplitude: physical
external-leg normalization and the rest of the process are still needed.

It is useful to remove the overall \(\ii\) from the full electron
propagator and define
\begin{equation}
 S_R(p)=\ii\mathcal S_R(p),\qquad
 \mathcal S_R^{-1}(p)=\slashed p-m-\Sigma_R(p).
\end{equation}
Gauge invariance then takes the particularly simple form
\begin{equation}
 \boxed{q_\mu\Gamma_R^\mu(p',p)
 =\mathcal S_R^{-1}(p')-\mathcal S_R^{-1}(p).}
 \label{eq:WT}
\end{equation}
At tree level this says \(\slashed q=\slashed p'-\slashed p\).
At one loop it relates the longitudinal vertex correction to the
momentum dependence of the self-energy computed in Lecture 28.

\begin{keyidea}[Longitudinal information does not determine the whole vertex]
The Ward--Takahashi identity constrains \(q_\mu\Gamma^\mu\).
A term proportional to \(\sigma^{\mu\nu}q_\nu\) is transverse because
\(\sigma^{\mu\nu}\) is antisymmetric. It is therefore unconstrained by
this contraction. That finite term carries the anomalous magnetic moment.
\end{keyidea}
The identity is off shell. Between external on-shell spinors, after
consistent leg normalization, its right side vanishes and it reduces to
the physical-photon Ward identity used for Compton scattering.
For loop calculations we keep momenta off shell when extracting
ultraviolet poles, so that soft singularities do not masquerade as
ultraviolet ones.

Throughout, \(\eta=(+---)\), \(\alpha=e^2/(4\pi)\),
\(d=4-2\eps\), and \(1/\bar\eps=1/\eps-\gamma_E+\ln4\pi\).
The symbol \(e\) is a signed charge. We use \(M\) for the physical
perturbative electron mass in on-shell form factors.

\module{Contact terms and charge renormalization}
Current conservation inside a time-ordered product is not simply zero:
derivatives also act on the ordering step functions. The equal-time
charge commutators are
\begin{equation}
 [Q,\psi]=-e\psi,\qquad [Q,\bar\psi]=e\bar\psi.
\end{equation}
Consequently, with the same renormalization prescription for the current,
\begin{align}
 \partial_\mu^x\langle\Tord j^\mu(x)\psi(y)\bar\psi(z)\rangle
 ={}&-e\delta^{(4)}(x-y)\langle\Tord\psi(y)\bar\psi(z)\rangle
 \nonumber\\
 &+e\delta^{(4)}(x-z)\langle\Tord\psi(y)\bar\psi(z)\rangle.
\end{align}
These endpoint terms are the operator counterpart of the diagrammatic
cancellations in Lecture 27. Fourier transforming and amputating the
external electron propagators gives \eqref{eq:WT}. A gauge-preserving
regulator and counterterms are understood; vector QED has no anomaly
that obstructs this identity.

Taking \(p'\to p\), with a regulator in place, gives
\begin{equation}
 \Gamma_R^\mu(p,p)=\gamma^\mu-\frac{\partial\Sigma_R(p)}{\partial p_\mu}.
 \label{eq:wardzero}
\end{equation}
Thus a change in the normalization of the propagator is accompanied by
a corresponding change in the zero-momentum vertex. The bare version
of this statement, expressed in renormalized fields, requires
\begin{equation}
 Z_1=Z_2,\qquad
 e_0=\mu^\eps e\frac{Z_1}{Z_2\sqrt{Z_3}}
     =\mu^\eps eZ_3^{-1/2}.
 \label{eq:Zrelation}
\end{equation}
The equality includes finite terms when a compatible normalization
scheme is used. The electron factors do not independently renormalize
the measured electric charge; their effect is tied together by the identity.

For example, the \(\MSbar\) ultraviolet self-energy is
\(\Sigma_{\rm UV}=a[-\xi\slashed p+(3+\xi)m]/\bar\eps\), where
\(a=\alpha/(4\pi)\). Differentiation predicts
\begin{equation}
 \Lambda^\mu_{\rm UV}=\frac{\xi a}{\bar\eps}\gamma^\mu,
 \qquad \delta_1=\delta_2=-\frac{\xi a}{\bar\eps}.
\end{equation}
The prediction provides a stringent check on the direct vertex integral.
It also shows why subtracting an arbitrary vertex pole without checking
the self-energy would be inconsistent with charge conservation.

\module{The one-loop integral and its ultraviolet part}
In Feynman gauge the vertex correction is
\begin{equation}
 \Lambda^\mu=-\ii e^2\mu^{2\eps}\int\frac{\dd^dk}{(2\pi)^d}
 \frac{\gamma^\alpha(\slashed p'-\slashed k+m)\gamma^\mu
 (\slashed p-\slashed k+m)\gamma_\alpha}
 {(k^2+\ii0)[(p'-k)^2-m^2+\ii0][(p-k)^2-m^2+\ii0]}.
 \label{eq:vertexloop}
\end{equation}
There is no closed-fermion-loop sign. For large \(k\), the numerator is
quadratic in \(k\) and the denominator scales as \(k^6\); the four-dimensional
integral is logarithmically divergent. The divergent part is independent
of external momenta and must be a local multiple of \(\gamma^\mu\).

To extract it, retain the leading numerator and use regulated angular
averaging:
\begin{equation}
 \gamma^\alpha\slashed k\gamma^\mu\slashed k\gamma_\alpha
 \ \longrightarrow\ 
 \frac{k^2}{d}\gamma^\alpha\gamma^\rho\gamma^\mu\gamma_\rho\gamma_\alpha
 =\frac{(d-2)^2}{d}k^2\gamma^\mu.
\end{equation}
After introducing Feynman parameters, the same term gives
\(\int_\ell\ell^2/(\ell^2-\Delta)^3\), whose ultraviolet pole is
\(\ii/(16\pi^2\bar\eps)\). The parameter triangle has area \(1/2\),
which cancels the factor two from combining three denominators. Thus
\begin{equation}
 \Lambda_{\rm UV}^\mu=\frac{\alpha}{4\pi\bar\eps}\gamma^\mu
 \quad(\xi=1).
\end{equation}
Do not set every scale to zero in the original integral: a scaleless
integral in dimensional regularization vanishes because ultraviolet and
infrared poles cancel, obscuring the ultraviolet coefficient being sought.

In a general covariant gauge, the longitudinal photon term adds
\((\xi-1)\alpha\gamma^\mu/(4\pi\bar\eps)\), agreeing with
\eqref{eq:wardzero}. In Landau gauge \(\xi=0\), this ultraviolet pole is
absent at one loop; the finite vertex correction is still present.

\begin{remarkbox}[What has and has not been made finite]
The local counterterm \(\delta_1\gamma^\mu\) removes the ultraviolet
pole. On-shell vertex functions at nonzero momentum transfer can still
contain infrared singularities. Ultraviolet subtraction and the removal
of a soft-photon regulator are different operations.
\end{remarkbox}
The remaining calculation extracts the finite transverse spin response.
It is simpler than evaluating every scalar coefficient of the general
off-shell vertex, which has many more tensor structures.

\module{On-shell form factors and the magnetic moment}
For \(p^2=p'^2=M^2\), Lorentz covariance, parity, and current conservation
reduce the electron current matrix element to two form factors:
\begin{equation}
 \bar u(p')\Gamma^\mu_{\rm phys}u(p)
 =\bar u(p')\left[\gamma^\mu F_1(q^2)
 +\frac{\ii\sigma^{\mu\nu}q_\nu}{2M}F_2(q^2)\right]u(p),
 \quad \sigma^{\mu\nu}=\frac{\ii}{2}[\gamma^\mu,\gamma^\nu].
 \label{eq:formfactors}
\end{equation}
Here \(\Gamma_{\rm phys}\) includes external-leg normalization, or is
evaluated in a scheme with unit regulated pole residue. An on-shell
charge condition imposes \(F_1(0)=1\). At tree level
\(F_1=1\) and \(F_2=0\).
A term \(q^\mu F_3\) is excluded by current conservation for generic
\(q^2\), with the smooth zero-transfer limit understood. Parity-violating
or electric-dipole structures are absent in the QED theory considered here.

The Gordon identity follows by inserting the on-shell Dirac equations
on the two sides of \(\gamma^\mu\):
\begin{equation}
 \bar u'\gamma^\mu u
 =\frac1{2M}\bar u'[(p'+p)^\mu+\ii\sigma^{\mu\nu}q_\nu]u.
\end{equation}
It separates a convection current from a spin current. The coefficient
of the latter in \eqref{eq:formfactors} is \(F_1+F_2\). In the slowly
varying magnetic-field limit, the spin interaction is
\begin{equation}
 H_{\rm spin}=-\bm\mu\cdot\bm B,\qquad
 \bm\mu=\frac{e}{2M}[F_1(0)+F_2(0)]\bm\sigma.
\end{equation}
Since \(\bm S=\bm\sigma/2\), define
\begin{equation}
 g=2[F_1(0)+F_2(0)],\qquad
 a_e\equiv\frac{g-2}{2}=F_2(0).
\end{equation}
For an electron \(e<0\), so the magnetic moment points opposite to the
spin even though \(g\) is positive in this convention.

\begin{keyidea}[Why the Pauli form factor is a prediction]
A local Pauli interaction would have dimension five. Minimal
renormalizable QED has no independent coefficient for it. The one-loop
\(F_2\) is ultraviolet finite and determined by the measured charge and
mass, rather than by a new ultraviolet counterterm.
\end{keyidea}
At order \(\alpha\), external-leg factors multiplying a Pauli term would
contribute only at order \(\alpha^2\), since that term vanishes at tree
level. The one-loop \(F_2\) can therefore be extracted directly from
\eqref{eq:vertexloop} between on-shell spinors.

\module{Isolating the finite Pauli coefficient}
Assign parameters \(x,y\) to the two electron denominators and \(z\) to
the photon denominator, with \(x+y+z=1\). Write
\begin{equation}
 \int_\triangle\equiv\int_0^1\dd x\,\dd y\,\dd z\,
 \delta(1-x-y-z),\qquad
 \frac1{ABC}=2\int_\triangle\frac1{(xA+yB+zC)^3}.
\end{equation}
Choose \(A=(p-k)^2-M^2\), \(B=(p'-k)^2-M^2\), \(C=k^2\), and shift
\(\ell=k-xp-yp'\). On shell,
\begin{equation}
 xA+yB+zC=\ell^2-\Delta,\qquad
 \Delta=M^2(1-z)^2-xyq^2.
 \label{eq:vertexdelta}
\end{equation}
This denominator retains the full momentum-transfer dependence.

Here is an economical way to organize the numerator algebra. Discard
terms odd in \(\ell\) under the regulated integral, use the external
Dirac equations only when slashed momenta reach the appropriate spinor,
and use the Gordon identity to eliminate \((p+p')^\mu\). Terms odd under
\(x\leftrightarrow y\) integrate to zero. For the finite Pauli coefficient
we may then set \(d=4\). The numerator becomes, in this integrated
on-shell sense,
\begin{align}
 N^\mu\doteq{}&
 [\ell^2-2M^2(1-4z+z^2)-2q^2(z+xy)]\gamma^\mu
 \nonumber\\
 &-4M^2z(1-z)\frac{\ii\sigma^{\mu\nu}q_\nu}{2M}.
 \label{eq:projectedN}
\end{align}
For checks on the reduction, use
\(\slashed q\gamma^\mu\slashed q=2q^\mu\slashed q-q^2\gamma^\mu\)
and \(\bar u'\slashed q u=0\). The angular average of
\(-2\slashed\ell\gamma^\mu\slashed\ell\) is \(\ell^2\gamma^\mu\).
The \(\doteq\) symbol is not an off-shell matrix identity.

The Pauli term has no \(\ell^2\) numerator and hence no ultraviolet pole.
Using
\begin{equation}
 \int\frac{\dd^4\ell}{(2\pi)^4}
 \frac1{(\ell^2-\Delta+\ii0)^3}
 =-\frac{\ii}{32\pi^2(\Delta-\ii0)},
\end{equation}
including the factor \(-2\ii e^2\) from the loop and parameterization,
gives
\begin{equation}
 \boxed{\displaystyle
 F_2(q^2)=\frac{\alpha}{\pi}\int_\triangle
 \frac{M^2z(1-z)}{M^2(1-z)^2-xyq^2-\ii0}+O(\alpha^2).}
 \label{eq:F2}
\end{equation}
The ultraviolet-divergent \(\gamma^\mu\) coefficient must instead be
computed with its full \(d\)-dependence; \eqref{eq:projectedN} is not a
shortcut for its finite part.

\module{Schwinger's result and its low-energy meaning}
At zero momentum transfer, the integrand in \eqref{eq:F2} depends only
on \(z\). At fixed \(z\), the allowed \(x\) interval has length \(1-z\).
Therefore
\begin{equation}
 F_2(0)=\frac{\alpha}{\pi}\int_0^1\dd z\,(1-z)
 \frac{z(1-z)}{(1-z)^2}
 =\frac{\alpha}{\pi}\int_0^1z\,\dd z
 =\frac{\alpha}{2\pi}.
\end{equation}
The apparent endpoint singularity of the triangle integrand is
integrable after its measure is included. Thus
\begin{equation}
 a_e=\frac{\alpha}{2\pi}+O(\alpha^2),\qquad
 g=2\left[1+\frac{\alpha}{2\pi}+O(\alpha^2)\right].
\end{equation}
This is the leading QED anomalous magnetic moment. It is independent of
the mass because the mass factors cancel in the dimensionless ratio;
the magnetic moment itself still scales as \(1/M\).

One may check infrared finiteness by introducing a small photon mass
\(\lambda\) solely as an auxiliary regulator in this finite integral.
Then \(\Delta\) becomes \(\Delta+z\lambda^2\), and
\begin{equation}
 F_2(0;\lambda)=\frac{\alpha}{\pi}\int_0^1\dd z\,
 \frac{z(1-z)^2}{(1-z)^2+z(\lambda/M)^2}
 \longrightarrow\frac{\alpha}{2\pi}.
\end{equation}
The integrand is bounded by \(z\), so its limit may be taken under the
integral. This auxiliary test does not replace the gauge-preserving
regulator used to establish the Ward identity and ultraviolet counterterms.

For slowly varying external fields, the result can be encoded by
\begin{equation}
 \Lag_{\rm Pauli}=-\frac{e\kappa}{4M}
 \bar\psi\sigma^{\mu\nu}\psi F_{\mu\nu},\qquad
 \kappa=F_2(0).
\end{equation}
With \(q=p'-p\), this produces the extra vertex
\(-\ii e[\ii\sigma^{\mu\nu}q_\nu\kappa/(2M)]\).
It is a low-energy matching coefficient of QED. At larger transfer the
full function \(F_2(q^2)\), together with \(F_1(q^2)\), is required.

\begin{remarkbox}[Finite magnetic response, infrared-sensitive charge response]
\(F_2(0)\) has a finite soft limit at this order. The ultraviolet-renormalized
\(F_1(q^2)\) at nonzero transfer generally retains soft-photon singularities.
The normalization \(F_1(0)=1\) is a charge condition; it does not assert
that the entire vertex is free of infrared effects.
\end{remarkbox}

\module{Soft photons and what an experiment measures}
A detector cannot distinguish a final state with no extra photon from
one accompanied by sufficiently low-energy radiation. The exclusive
Fock-state probability for exactly no soft radiation is therefore not
the quantity to compare with a finite-resolution measurement.

The leading soft-emission factor can be derived from an external
fermion line. If a photon of momentum \(k\) is emitted from an outgoing
line of momentum \(p\), its adjacent propagator has
\((p+k)^2-M^2\simeq2p\cdot k\). The Dirac equation reduces the numerator
to \(2p^\mu\). Summing attachments to all charged external lines gives
\begin{equation}
 \cM_{n+\gamma}(k,\epsilon)
 \simeq\left[\sum_i\eta_i e_i
 \frac{p_i\cdot\epsilon^*}{p_i\cdot k}\right]\cM_n,
 \quad
 \eta_i=\begin{cases}+1&\text{outgoing},\\-1&\text{incoming}.
 \end{cases}
\end{equation}
Here \(e_i\) is the signed physical charge of each particle, including
antiparticles. A gauge change \(\epsilon\mapsto\epsilon+ck\) gives a
factor \(c\sum_i\eta_i e_i=0\), precisely charge conservation.

For massive external particles, integrating unresolved real photons
with \(\omega<\Delta E\) produces the factor
\begin{equation}
 \int_{\omega<\Delta E}\frac{\dd^3k}{(2\pi)^3\,2\omega}
 \sum_{\lambda=1}^2
 \left|\sum_i\eta_i e_i
 \frac{p_i\cdot\epsilon_\lambda^*(k)}{p_i\cdot k}\right|^2.
\end{equation}
The soft behavior is \(\int\dd\omega/\omega\). A virtual photon
connecting charged lines gives the same infrared logarithm with the
opposite contribution to the probability. Consequently the inclusive
combination
\begin{equation}
 \dd\sigma_{\rm incl}(\Delta E)
 =\dd\sigma_{\rm Born}+\dd\sigma_{\rm virtual}
 +\dd\sigma_{\rm real,\,\omega<\Delta E}
\end{equation}
is finite as the soft regulator is removed, to the stated perturbative
order. Physical logarithms of the resolution \(\Delta E\) remain.

The cancellation requires the complete virtual correction, including
external legs and other diagrams of the process, not the vertex alone.
We assumed nonzero charged-particle masses; a massless limit introduces
additional collinear singularities and requires a suitable inclusive
observable. These statements explain the infrared qualification to LSZ
in Lectures 23 and 28 while preserving the usefulness of perturbative
amplitudes for measured rates.

\module{The QED renormalization pattern}
Collect the one-loop ultraviolet constants for one Dirac species in a
covariant gauge. With \(a=\alpha/(4\pi)\),
\begin{equation}
 Z_3=1-\frac{4a}{3\bar\eps},\quad
 Z_2=Z_1=1-\frac{\xi a}{\bar\eps},\quad
 Z_m=1-\frac{3a}{\bar\eps}.
\end{equation}
These are \(\MSbar\) constants, not physical pole residues. All ultraviolet
poles in the two-point functions and vertex are absorbed by the original
QED operators. The equality \(Z_1=Z_2\) is the consistency condition
linking the calculations, and a photon mass is not generated.

One can read off the running charge directly. From \eqref{eq:Zrelation},
\begin{equation}
 e_0=\mu^\eps e\left[1+\frac{e^2}{24\pi^2\bar\eps}+O(e^4)\right].
\end{equation}
Hold \(e_0\) fixed and differentiate with respect to \(\ln\mu\).
The leading \(d\)-dimensional running is \(-\eps e\); retaining its
product with the pole term gives
\begin{equation}
 \beta_e=-\eps e+\frac{e^3}{12\pi^2}+O(e^5),\qquad
 \beta_\alpha\big|_{d=4}=\frac{2\alpha^2}{3\pi}+O(\alpha^3).
\end{equation}
This agrees with the logarithm of the polarization function in Lecture
28. For multiple active species, multiply the one-loop coefficient by
\(\sum_f Q_f^2\), where their charges are \(Q_f e\); threshold matching
is needed when a species becomes heavy compared with the process scale.

\begin{center}
\begin{tabular}{p{0.27\textwidth}p{0.61\textwidth}}
\toprule
Quantity & Main physical information \\
\midrule
Transverse photon two-point function & Charge screening, running, and pair thresholds \\
Electron self-energy & Mass relation and regulated external-leg normalization \\
Longitudinal vertex identity & Equality of vertex and electron field renormalization \\
Pauli form factor & Finite spin response, \(a_e=\alpha/(2\pi)+\cdots\) \\
Inclusive probability & Cancellation of unobservable soft-regulator dependence \\
\bottomrule
\end{tabular}
\end{center}
Renormalization has not removed the momentum dependence of QED. It has
separated unphysical ultraviolet parameters from finite predictions
expressed in measured quantities. Gauge identities coordinate that
separation, while unitarity and inclusiveness determine how the resulting
amplitudes enter observable probabilities.

Homework 9 follows the same chain: photon constraints and polarizations,
a gauge-invariant scattering amplitude, propagator corrections, and the
vertex identity with the magnetic moment. Together these are the basic
tools needed for systematic perturbative calculations in QED.
\end{document}
